% !PRINTED 493
\weberpage{493}

Wir beweisen noch den Satz:

5. \gesperrt{Jede algebraische Zahl $\Theta$ l\"asst sich durch Multiplication mit
einer nat\"urlichen Zahl in eine ganze algebraische Zahl verwandeln.}

Denn ist
\[
\Theta^m + A_1 \Theta^{m-1} + \cdots + A_m = 0,
\]
und sind $A_1, \ldots, A_m$ rationale Zahlen mit dem gemeinsamen Nenner $a$, so
erh\"alt man durch Multiplication mit $a^m$:
\[
(a\Theta)^m + A_1 a (a\Theta)^{m-1} + A_2 a^2 (a\Theta)^{m-2} + \cdots + A_m a^m = 0,
\]
woraus hervorgeht, dass $a\Theta$ eine \gesperrt{ganze Zahl} ist.

\begin{center}
\S.\,134.

\gesperrt{Algebraische K\"orper.}
\end{center}

Im dreizehnten Abschnitte des ersten Bandes haben wir gesehen, wie man aus
jeder Wurzel $\Theta$ einer in irgend einem K\"orper $\Omega$ irreduciblen
Gleichung $n^{\mathrm{ten}}$ Grades $f(\Theta) = 0$ einen algebraischen
K\"orper $\Omega(\Theta)$ \"uber $\Omega$ ableitet. Die aus den $n$ Wurzeln
dieser Gleichung abgeleiteten $n$ K\"orper, die auch zum Theil oder alle
identisch sein k\"onnen, haben wir \gesperrt{conjugirte K\"orper} genannt.

Bezeichnen wir mit $R$ den K\"orper der rationalen Zahlen, so giebt also nach
unserer Definition jede algebraische Zahl $n^{\mathrm{ten}}$ Grades, $\Theta$,
Anlass zu einem algebraischen K\"orper $R(\Theta)$ \"uber $R$, den wir von
jetzt an kurz einen algebraischen Zahlk\"orper $n^{\mathrm{ten}}$
Grades nennen. Wir haben auch schon fr\"uher nachgewiesen (Bd.\,I,
\S.\,143), dass man immer einen \gesperrt{algebraischen Zahlk\"orper bestimmen kann,
der eine endliche Anzahl beliebig gegebener algebraischer Zahlen enth\"alt.}

Dieser Satz wird an dieser Stelle hervorgehoben, um darauf hinzuweisen, dass
die Allgemeinheit einer Betrachtung \"uber irgend eine endliche Anzahl
algebraischer Zahlen dadurch nicht beeintr\"achtigt wird, dass man diese
Zahlen alle in einem algebraischen Zahlk\"orper gelegen voraussetzt.

Jede Zahl $\omega$ eines solchen K\"orpers kann als ganze Function
$\varphi(\Theta)$ von $\Theta$ mit rationalen Co\"efficienten dargestellt
werden, und jeder Zahl $\omega$ entspricht in jedem der $n$ conjugirten
K\"orper eine bestimmte Zahl. Diese \gesperrt{conjugirten Zahlen} k\"onnen zum Theil
einander gleich sein, und wir haben danach primitive und
%CONT>>
